\documentclass[11pt]{clean_cv}

\usepackage[T1]{fontenc}
\usepackage{lmodern}
\usepackage{pacioli}
\usepackage{bookman}
\usepackage{amsmath}

\renewcommand{\baselinestretch}{0.9}

\addbibresource{publications.bib}



\author{
{\boldshape
L\hspace{1px}U\hspace{1px}I\hspace{1px}G\hspace{1px}I\hspace{20px}P\hspace{1px}I\hspace{1px}Z\hspace{1px}Z\hspace{1px}A }
{}
}
\headlineposition{College Student}

\begin{document}
\maketitle
% In this section, you can use any of the FontAwesome icons. The commands \faCenter and \faCenterStyle have been defined to properly center the icons
% when using the default font settings.
\vspace{2px}
\begin{center}
\begin{tabular}{llll}
% \href{tel:+393668341405}{\faCenter{phone-alt} +39 3668341405} &
\href{mailto:luigi.pizza2002@gmail.com}{\faCenter{envelope} luigi.pizza2002@gmail.com} & 
\href{https://www.linkedin.com/in/luigipizza/}{\faCenter{linkedin} Luigi Pizza} &
\href{https://github.com/luigi-pizza}{\faCenter{github} luigi-pizza}\\
\end{tabular}
\begin{flushright}
% \textbf{\small \normalfont Available to start in September/October 2026}
\end{flushright}
\end{center}

% \vspace*{-5px}

\vspace{-1.5em}




\vspace*{-20px}
\section{{\normalfont Education}}

\begin{datetabular}{3.2em}

\dateentry{Ongoing Sep 2024}{
\textbf{ETH Z\"urich} 
\hfill { Expected Graduation date: September 2026.} \hspace*{5px}

{\itshape Master’s Degree in  \href{https://inf.ethz.ch/studies/master/master-cs-2020.html}{Computer Science}. Majoring  in Machine Intelligence.}
}


\dateentry{Jul 2024 Sep 2021}{
\textbf{La Sapienza -- University of Rome} 
\hfill
Final Grade: {\fontfamily{lmr}\selectfont110/110} cum laude. \hspace*{5px}

\textit{Bachelor’s Degree in \href{https://corsidilaurea.uniroma1.it/en/corso/2022/30786/home}{Applied Computer Science and Artificial Intelligence}},
\hfill
GPA: {\fontfamily{lmr}\selectfont 4.0} (Avg. {\fontfamily{lmr}\selectfont29.89/30}), Top {\fontfamily{lmr}\selectfont 1.} \hspace*{5px}

Selected as the top one student in my cohort to take part in the \href{https://www.uniroma1.it/en/pagina/honours-programme}{Honours Program}.
}

\end{datetabular}

\vspace*{-15px}
\section{{\normalfont Work and Research Experience}}

\begin{datetabular}{3.2em}

\dateentry{Jul 2026 Sep 2025}{
  \textbf{
  \href{https://www.vvz.ethz.ch/Vorlesungsverzeichnis/lerneinheit.view?semkez=2025W&ansicht=ALLE&lerneinheitId=192925&lang=en}{Head Teaching Assistant -- Advanced Machine Learning}
  }
  \hfill ETH Zürich

  \textbf{
  \href{https://vvz.ethz.ch/Vorlesungsverzeichnis/lerneinheit.view?semkez=2025S&ansicht=ALLE&lerneinheitId=186921&lang=en}{Teaching Assistant --  Stochastics and Machine Learning}
  } 
   
    % \begin{itemize}[leftmargin=1.3em, labelsep=0.3em]
    %     \item Topics Covered: Gaussian processes, neural networks, NLP, transformers, diffusion models, CV, RL, and PAC learning. 
    % \end{itemize}
  
  \hspace{0.5em}{\tiny $\raisebox{1px}{\bullet}$}\hspace{0.2em}  Topics Covered: Gaussian processes, neural networks, NLP, transformers, diffusion models, CV, RL, and PAC learning. 
  
\hspace{0.5em}{\tiny $\raisebox{1px}{\bullet}$}\hspace{0.2em} Duties: Creation of lecture notes/material/assignments; conducted exercise sessions.
}



\dateentry{Dec 2026 Jan 2025}{
\textbf{\href{https://viscon.ethz.ch/2026/contact}{President, Head of Organization - VIScon2026}, \quad \href{https://viscon.ethz.ch/2025/contact}{Head of Symposium - VIScon2025}}

Organized multiple times one of the largest computer science symposiums in the Zürich area, featuring invited speakers presenting innovative ideas. The event attracts every year over {\fontfamily{lmr}\selectfont 
1,000} participants and includes a Hackathon and the Symposium.
}

% \dateentry{2025 2024 2023}{
% \textcolor{gray}{Participation to} \textbf{Summer Schools} 


% I have participated to \textbf{\href{https://cmmrs.mpi-sws.org/}{CMMRS 2025}}, \textbf{\href{https://boost24.elicsir.it/home/}{Boost 24}}, \textbf{\href{https://fa23.bici.events/}{FA 23}} and \textbf{\href{https://buca22.bici.events/home}{BUCA’22:}} \textit{Challenges in building Billion Users Cloud Applications}. In these summer schools, outstanding students are invited to learn about cutting-edge research in computer science from leading researchers and industry professionals.
% }



% % CMMRS has united top pre-doctoral students with faculty from Cornell, UMD, and Max Planck Institutes for a week-long school featuring cutting-edge research, networking, and mentorship. Participants gain insight into academic and industrial research careers and gain a head start for graduate school applications.
% }


% \dateentry{2024}{
% \textbf{\href{https://swerc.eu/2023/about/}{Southwestern Europe Regional Contest (SWERC) 2023-2024}}

% I've represented Sapienza (team \href{https://swerc.eu/2023/theme/scoreboard/domjudge.lip6.fr/public/team/48.html}{SapienzaRed}) at the \href{https://swerc.eu/2023/theme/scoreboard/domjudge.lip6.fr/public.html}{SWERC competition}: a programming contest for teams of three students, focused on algorithmic problem solving and practical coding. Placed $40^{th}$ out of $107$ teams.
% }

% \dateentry{2024}{
% \textbf{\href{https://boost24.elicsir.it/home/}{Boost'24:}}
% \textit{Bologna Orthogonal Summer Term}

% A summer-school where outstanding students are invited to learn about cutting-edge research in computer science from leading researchers.
% }






% \dateentry{2023}{
% \textbf{\href{https://www.uniroma1.it/en/pagina/honours-programmes}{Honours Program}}

% The Honours Program is integrated with the regular study plan and is addressed to students willing to deepen their knowledge and capabilities through extracurricular activities in selected topics. My selected Area of Interest is "Blockchain technology" under the supervision of Professor Alessandro Mei.
% }









% \begin{datetabular}{4.2em}

% \dateentry{2023}{
% \textbf{\href{https://fa23.bici.events/}{FA'23:}} 
% \textit{il Futuro Annunciato}

% A brief workshop aiming to provide insight into the new frontiers of computer science through seminars, interactive panel discussions, and informal conversations with computer researchers from academia, industry leaders, and startup founders.
% }

% % \dateentry{2022}{
% % \textbf{\href{https://buca22.bici.events/home}{BUCA’22:}} \textit{Challenges in building Billion Users Cloud Applications}

% % Summer School revolving around the challenges of building billion-user systems and ways to address these challenges. The Instructors are 3 High-Ranking Google Engineers specialised in this field. 
% % }




% % There is something seriously funky happening between the tabular commands and the end of the itemize blocks.
% % The command \eatvspace should be used if extra space appears at the end of a datetabular environment.


\end{datetabular}

\vspace*{-15px}
\section{{\normalfont Projects}}

\begin{datetabular}{3.2em}

\dateentry{2025}{
  \textbf{\href{https://github.com/rubenciranni/cuda-bopm}{American Option Pricing CUDA}}, ETH Zürich 
  \hfill Final Grade: 6/6 

    Designed and implemented (CUDA, C++) in a team of five a high-performance implementation of the binomial model for pricing American options.
}

\dateentry{2025}{
  \textbf{\href{https://doi.org/10.1145/3765964.3811657}{Semester Project – Diffusion Models}}, ETH Zürich 
  \hfill Ongoing Progect 

    Conducting research on effective teaching methods for diffusion processes in machine learning, with the goal of improving conceptual understanding and intuition behind diffusion-based generative models. Resulted in publication in 
}


\dateentry{2025}{
  \textbf{Affordance Understanding in 3D Environments}, ETH Zürich 
    \hfill Final Grade: {\fontfamily{lmr}\selectfont 5.75/6}
    
    Designed and implemented (Python + PyTorch) in a team of four a novel approach for open-vocabulary 3D segmentation
    in indoor environments. The pipeline can segment functional elements (e.g., a door handle) in a 3D point cloud of an
    indoor environment given a natural language prompt describing an action (e.g., "open the bedroom door").
}

\dateentry{2023}{
\textbf{\href{https://github.com/luigi-pizza/SelfDrivingDrone}{SelfDrivingDrone}}, Sapienza 
\hfill Final Grade: {\fontfamily{lmr}\selectfont 30/30} cum laude

Designed and developed (Python + PyTorch) in a team of three, this project develops an autonomous obstacle navigation system for the DJI Tello drone. The drone is capable of autonomous flight through a predetermined set of obstacles. 
}



\end{datetabular}
\vspace*{-15px}
\section{{\normalfont Additional Experience and Awards}}

\begin{datetabular}{3.2em}

% % CMMRS has united top pre-doctoral students with faculty from Cornell, UMD, and Max Planck Institutes for a week-long school featuring cutting-edge research, networking, and mentorship. Participants gain insight into academic and industrial research careers and gain a head start for graduate school applications.
% }

\dateentry{2024}{
\textbf{\href{https://swerc.eu/2023/about/}{Southwestern Europe Regional Contest (SWERC) 2023-2024}}

\textit{A programming contest for teams of three students, focused on algorithmic problem solving and practical coding.}

I've represented Sapienza (team \href{https://swerc.eu/2023/theme/scoreboard/domjudge.lip6.fr/public/team/48.html}{SapienzaRed}) at the \href{https://swerc.eu/2023/theme/scoreboard/domjudge.lip6.fr/public.html}{SWERC competition}: a  Placed $40^{th}$ out of $107$ teams.
}


\dateentry{2023}{
\textbf{\href{https://cyberchallenge.it/}{CyberChallenge.IT}}

\textit{Participated in the main Italian initiative to identify and recruit the next generation of CyberSecurity professionals.}
}


\dateentry{Sep 2024 Sep 2021}{
\textbf{\href{https://www.collegiocavalieri.it/}{Residence Hall of the Cavalieri del Lavoro "Lamaro Pozzani"}} 

{ \itshape Residence in a merit-based institution recognized by the Italian Ministry of University and research.} 

Attended here courses on business culture, and various events by the italian business élite. }


\dateentry{2025 2022}{
\textcolor{gray}{Participation to} \textbf{Summer Schools} and \textbf{Civic Engagement Workshops} 

\hspace{0.5em}{\tiny $\raisebox{1px}{\bullet}$}\hspace{0.2em}  Sat in a panel discussion at the European Parliament to discuss \href{https://www.linkedin.com/posts/luigipizza_eye2023-strasbourg-activecitizenship-activity-7140668305626902528-aroh/}{"The Role of Student Affairs \& Services in Skills Development".}

\hspace{0.5em}{\tiny $\raisebox{1px}{\bullet}$}\hspace{0.2em}  Participation to various summer schools for outstanding students: \textbf{\href{https://cmmrs.mpi-sws.org/}{CMMRS 2025}}, \textbf{\href{https://boost24.elicsir.it/home/}{Boost 24}}, \textbf{\href{https://fa23.bici.events/}{FA 23}} and \textbf{\href{https://buca22.bici.events/home}{BUCA’22}}. 
% In these summer schools, outstanding students are invited to learn about cutting-edge research in computer science from leading researchers and industry professionals.

\hspace{0.5em}{\tiny $\raisebox{1px}{\bullet}$}\hspace{0.2em}
Participation to projects sponsored by the European Parliament: \textit{\href{https://www.linkedin.com/posts/luigipizza_eye2023-strasbourg-activecitizenship-activity-7140668305626902528-aroh}{EYE2023 Tournament}, \href{https://www.euchangemakers.com/genz-dubrovnik-croatia}{GenZ Dubrovnik}, \href{https://www.euchangemakers.com/egdszeged}{EGD Szeged}, \href{https://www.euchangemakers.com/egdrome}{EGD Rome}.}
}

% \dateentry{2024 2023 2022}{
% \textcolor{gray}{Participation to} \textbf{\href{https://www.euca.eu/}{EucA} Projects} \quad \textit{\href{https://www.linkedin.com/posts/luigipizza_eye2023-strasbourg-activecitizenship-activity-7140668305626902528-aroh}{EYE2023 Quiz Tournament}, \href{https://www.euchangemakers.com/genz-dubrovnik-croatia}{GenZ Dubrovnik}, \href{https://www.euchangemakers.com/egdszeged}{EGD Szeged}, \href{https://www.euchangemakers.com/egdrome}{EGD Rome}}

% Projects by \href{https://www.euca.eu/}{EucA} and supported by the European Parliament, tackling civic participation, guerrilla marketing and cyber security. Through EucA, of which I am an ambassador, I had the opportunity to \href{https://www.linkedin.com/posts/luigipizza_eye2023-strasbourg-activecitizenship-activity-7140668305626902528-aroh/}{sit in a panel discussion at the European Parliament to speak of "The Role of Student Affairs \& Services in Skills Development".}
% }

\dateentry{2021}{
\href{http://olimpiadi.dm.unibo.it/2021/05/28/finale-nazionale-della-gara-individuale-tabellone-2021/}{Bronze Medal in the Italian Maths Olympiads}
}


\end{datetabular}
\vspace*{-15px}
\section{{\normalfont Skills}}

\begin{tabbing}
\qquad
\= \small \textbf{Parallel and GPU Programming}: \qquad \= MPI, PThreads, CUDA \\
\> \small \textbf{Programming Languages}: \> \small Python, C, C++, Go, Java, R \\
\> \small \textbf{Machine Learning}: \> \small PyTorch, Lightning, Scikit-learn, Webscraping, Selenium \\
\> \small \textbf{Data Analytics}: \> \small pandas, PyTorch, NumPy, Matplotlib, Seaborn, Plotly, Scikit-learn, R \\
\> \small \textbf{Databases}: \> \small SQL, PostgreSQL \\
\> \small \textbf{Tools}: \> \small Git, Anaconda, VS Code, Docker, Algorithms \& Data Structures \\
\> \small \textbf{Languages}: \> \small English, Italian \\
\end{tabbing}

% \dateentry{\color{black} \small
% Project Management Tools:

% Programming Languages:

% Machine Learning: 

% Data Analytics:

% Databases:

% Languages:}
% {
% Git, Anaconda

% C, C++, Java, Python

% Scikit-learn, Lightning, PyTorch

% pandas, scikit-learn, R

% SQL

% English, Italian
% }
% \end{datetabular}

% \section{Languages}

% \begin{center}\\[7px]
% \begin{tabular}{lllll}
% \textbf{Italian}  & Native Speaker &  &
% \textbf{English}  & C1-C2           \\
% %\textbf{Spanish}  & A2-B1 & & \textbf{French}  & A2-B1
% \\
% \end{tabular}
% \end{center}


\eatvspace

\begin{datetabular}{5em}
\dateentry{Dec 2023}{
\textbf{IELTS Academic Overall Band Score 8}
% Obtained IELTS English certification
, with CEFR level C1/C2.
}

\end{datetabular}



% \section{Projects}


% \begin{datetabular}{2em}


% \dateentry{2023}{
% \textbf{\href{https://github.com/luigi-pizza/ParallelPCA}{ParallelPCA}} 
% \textit{Multicore Programming}

% Designed and developed with a team of two other students, this project implements a distributed algorithm for Principal Components Analysis. ParallelPCA implements two different methodologies: multithreading and message passing between processes. 
% }


% \dateentry{2023}{
% \textbf{\href{https://github.com/luigi-pizza/wasaPhoto}{WASAPhoto}} 
% \textit{Full-Stack Web Application}

% This project develops a social media platform that focuses on photo sharing. It implements a Go backend and Vue.js frontend.
% }


% \dateentry{2023}{
% \textbf{\href{https://github.com/luigi-pizza/SelfDrivingDrone}{SelfDrivingDrone}} 
% \textit{Computer Vision}

% Designed and developed with a team of 2 other students, This project develops an autonomous obstacle navigation system for the DJI Tello drone. The drone is capable of autonomous flight through a predetermined set of obstacles. 
% }



% \dateentry{2022}{
% \textbf{\href{https://github.com/luigi-pizza/BattleOfSexesSimulation}{Population simulator}} \textit{Java, JavaFX}

% Designed and developed with a team of 2 other students, an evolution simulator in Java based on the 1976 book "The Selfish Gene" by Richard Dawkins, using multithreading to achieve a non-deterministic result. A GUI implemented with JavaFX enables the visualization of the simulation in real-time.
% }


% \end{datetabular}



% \section{Skills}
% \begin{datetabular}{15em}
% \dateentry{
% Programming Languages:

% Mathematics:

% Computer Architecture:

% AI: 

% Databases:

% Project Management Tools:}
% {C, C++, Java, Python

% Multivariate Calculus, Linear Algebra, Probability, Statistics

% x86, RISC-V, Operative Systems, Networks

% Scikit-learn, PyTorch, Computer Vision, NLP

% SQL, MongoDB, Neo4J

% Git, Anaconda
% }
% \end{datetabular}


%\centerline{\textbf{Italian}  \quad Native Speaker \quad\quad\quad\quad\quad\quad\quad\quad
%\quad\quad\quad\quad\quad\quad\quad\quad\textbf{English}  \quad C1-C2}
%\\
%\quad\quad\quad\quad\quad\quad\quad\quad\quad\quad\textbf{Spanish}  \quad A2-B1 \quad\quad\quad\quad\quad\quad\quad\quad
%\quad\quad\quad\quad\quad\quad\quad\quad\quad\quad\textbf{French}  \quad A2-B1
%\\
%\section{Publications}
%\nocite{*} % Loads every entry from the attached .bib file
% Highlight takes three entries, given name, given name initials, and family name.
% If you have a middle initial, this call looks like:
% \highlightauthorname{Bob H.}{B. H.}{Smith}
%\highlightauthorname{Frodo}{F.}{Baggins} 
%\begin{datetabular}{2em}
%printbibyear has been defined to only print entries from a given citation year.
%\dateentry{3011}{\printbibyear{3001}}
%\dateentry{3001}{\printbibyear{3011}}
%\end{datetabular}

\end{document}
